\section{Strict vanishing through the coefficient-prime character tower}
\label{sec:cubic-prime-character-strict-vanishing}

The two residual calculations over the original cubic base control
more than two individual Dirichlet characters.  The strict local
condition at the coefficient prime and the unchanged away inertia
allow every lift in the following character families.  We give the
coefficient argument explicitly; it uses no automorphic host or
control theorem.

\paragraph{Characters and original coefficients.}
Keep \(p=269\), \(K=\mathbb Q(\sqrt{-3})\), either generator
\(\chi\) of the cubic character of \(K(\sqrt[3]{269})/K\),
and \(j\in\{1,2\}\).  Let \(c_p\) be the actual action on
roots of unity, and let
\[
 \omega_p:G_{\mathbb Q}\longrightarrow\mu_{268}\subset\mathbb Z_p^\times
 \quad\text{be the Teichmuller lift of }\bar c_p.
 \tag{G125.1}
\]
This convention uses the action on roots of unity; in particular
geometric Frobenius has cyclotomic value \(q^{-1}\).
The lift exists uniquely because each nonzero residue class has a
unique root of \(X^{268}-1\) above it, by the simple-root form of
Hensel's lemma.  Uniqueness makes the lift multiplicative.
Put \(\psi_0=\chi_{-4}|_{G_K}\).

Let \(E/\mathbb Q_p\) be any finite extension containing the
cubic values, with its \emph{original} valuation ring \(O\),
uniformizer \(\varpi\), maximal ideal \(\mathfrak m\), and
residue field \(k\).  Fix a continuous character
\[
 \kappa:G_K\longrightarrow 1+\mathfrak m,
 \qquad \kappa|_{I_v}=1\quad(v\nmid p).
 \tag{G125.2}
\]
No extension of \(\kappa\) to \(G_{\mathbb Q}\), finite-order
hypothesis, or parity hypothesis is imposed.  For integers
\(e\ge1\) and \(r\), define
\[
 a_{j,e,r,\kappa}
   =\chi^{-j}\psi_0^{-1}\omega_p^{-r}\kappa^{-1}c_p^{-e}|_{G_K},
 \qquad D_{j,e,r,\kappa}=(E/O)(a_{j,e,r,\kappa}).
 \tag{G125.3}
\]
The notation for restriction in this formula applies to the
rational factors; \(\chi,\psi_0,\kappa\) are already characters
of \(G_K\).  All actions on \(E/O\) are continuous for its
discrete topology: the stabilizer of an element killed by
\(\varpi^n\) contains the open kernel of the character modulo
\(\varpi^n\).

Use exactly the original strict local kernel
\[
 \operatorname{Sel}_{\rm str}(K,D)=\ker\left(
 H^1_{\rm cts}(G_K,D)\longrightarrow
 \prod_{v\mid p}H^1_{\rm cts}(G_{K_v},D)
 \times\prod_{v\nmid p,\infty}H^1_{\rm cts}(I_v,D)\right).
 \tag{G125.4}
\]
The coefficient-prime condition is full local zero.  No finite
Bloch--Kato condition is substituted for it.

\begin{theorem}[AC-G125-lifts: every original continuous lift in the two branches]
\label{thm:ACG125lifts}
Under (G125.1)--(G125.4), if
\[
                         e+r\equiv133\text{ or }267\pmod{268},
 \tag{G125.5}
\]
then, for both cubic labels,
\[
 \operatorname{Sel}_{\rm str}(K,D_{j,e,r,\kappa})=0.
 \tag{G125.6}
\]
Its length and normalized length over that same original \(O\)
are consequently zero.  There is no condition
\(\operatorname{length}_O(O/p)=1\).
The only finite arithmetic inputs are the residual vanishing
statements AC-G114-residual and AC-G115-residual, with their
previously certified class and explicit unit calculations.
\end{theorem}

\begin{proof}
\emph{1. Identify the unchanged residual representation.}
Reduction of (G125.3) gives
\[
 \bar a_{j,e,r,\kappa}
       =\bar\chi^{-j}\bar\psi_0^{-1}\bar c_p^{-(e+r)},
 \tag{G125.7}
\]
because \(\bar\kappa=1\) and \(\bar\omega_p=\bar c_p\).
Choose \(e_0\in\{133,267\}\) congruent to \(e+r\).
Equation (G125.7) is precisely the coefficient character of
AC-G115-residual if \(e_0=133\), and of AC-G114-residual
if \(e_0=267\).  Those results hold over every finite residue
field containing the character values.  Therefore
\[
       \operatorname{Sel}_{\rm str}(K,k(\bar a_{j,e,r,\kappa}))=0.
 \tag{G125.8}
\]
This identifies actual \(G_K\)-modules and their full local
restriction maps.  Equality of residual dimensions alone would
not suffice, and is not the argument used here.

\emph{2. Check the invariant modules for this full lift.}
There is one prime of \(K\) over \(269\); the cubic extension
there is totally ramified of degree three.  Choose an inertia
element with nontrivial \(\bar\chi\)-value and raise it to the
\(268\)th power.  Its rational residual cyclotomic value becomes
one, while its cubic value remains nontrivial since
\(268\equiv1\pmod3\).  The character \(\psi_0\) is unramified
at this prime and \(\bar\kappa=1\).  Thus there is an element
\(g\) of coefficient-prime inertia for which
\(\bar a_{j,e,r,\kappa}(g)\ne1\).
Consequently \(a_{j,e,r,\kappa}(g)-1\) is a unit of \(O\).
Its multiplication on \(E/O\) is invertible, proving
\[
                  D^{G_{K_p}}=D^{G_K}=0,
             \qquad D=D_{j,e,r,\kappa}.
 \tag{G125.9}
\]
This check is residual and remains valid even if \(\kappa\)
has infinite image or the original coefficient extension is ramified.

At every finite \(v\nmid p\), the characters \(c_p\),
\(\omega_p\), and \(\kappa\) are trivial on inertia.
Hence
\[
                  a_{j,e,r,\kappa}|_{I_v}
                              =\chi^{-j}\psi_0^{-1}|_{I_v},
 \tag{G125.10}
\]
a character of order dividing six.  If it is trivial,
\(D^{I_v}=D\), which is \(\varpi\)-divisible.  If it is
nontrivial, some value is a nontrivial root of unity of order
prime to \(p\); its difference from one is a unit.  Indeed
reduction is injective on roots of unity of order prime to
\(p\), by the same simple-root argument used in (G125.1).
Then \(D^{I_v}=0\).  Thus every local invariant module used
in (G125.4) is \(\varpi\)-divisible, including those at
places ramified in \(\chi\) or \(\psi_0\).
The assumption that \(\kappa\) is unramified away from
\(p\) is used here; arbitrary new away ramification is not
permitted by this theorem.

\emph{3. Lift a strict uniformizer-torsion class without changing \(O\).}
The actual uniformizer gives the exact coefficient sequence
\[
 0\longrightarrow k(\bar a_{j,e,r,\kappa})
  \xrightarrow{\ \bar u\mapsto u/\varpi+O\ }D
  \xrightarrow{\varpi}D\longrightarrow0.
 \tag{G125.11}
\]
To check its cohomological use directly, let a continuous
cocycle \(z\) represent a strict class killed by \(\varpi\).
Write \(\varpi z=dm\) and choose \(n\in D\) with
\(\varpi n=m\).  Then \(w=z-dn\) takes values in
\(D[\varpi]\), identified with the first term of (G125.11).
If a residual cocycle becomes \(du\) in \(D\), then
\(\varpi u\) is invariant.  Globally it is zero by
(G125.9), so \(u\in D[\varpi]\).  Thus the residual
lift of the global class is unique.

For any local group in (G125.4), strictness makes the local
restriction of \(w\) a coboundary \(du_v\) in \(D\).
Its vanishing under \(\varpi\) makes \(\varpi u_v\)
invariant.  Divisibility of that local invariant module
provides invariant \(u_{v,0}\) with
\(\varpi u_{v,0}=\varpi u_v\).  The element
\(u_v-u_{v,0}\) then lies in \(D[\varpi]\) and has the
same coboundary.  Thus the residual lift is locally zero
in every prescribed group: it is a strict residual class.
By (G125.8) it is zero, proving
\[
                 \operatorname{Sel}_{\rm str}(K,D)[\varpi]=0.
 \tag{G125.12}
\]
All coboundaries here are continuous; every chosen coefficient
element has an open stabilizer as checked after (G125.3).

\emph{4. Eliminate every remaining primary-torsion class.}
A continuous map from compact \(G_K\) to discrete \(D\)
has finite image.  Each element of \(E/O\) is killed by
some power of the original uniformizer.  One common power
therefore kills all values of any continuous cocycle.
Every cohomology class, including every strict class, is
\(\varpi\)-primary torsion.  A nonzero strict class with
least annihilating power \(\varpi^n\) would have a
nonzero \(\varpi^{n-1}\) multiple in (G125.12).
This contradiction proves (G125.6).
\end{proof}

\begin{theorem}[AC-G125-Dirichlet: exact finite-character family and count]
\label{thm:ACG125Dirichlet}
Let \(m\ge1\), and view finite characters modulo \(269^m\)
as characters valued in \(\overline{\mathbb Q}_p^\times\) of
\(\operatorname{Gal}(\mathbb Q(\mu_{269^m})/\mathbb Q)\)
using the action-on-roots convention.  For any fixed \(e\ge1\),
exactly
\[
                              2\cdot269^{m-1}
 \tag{G125.13}
\]
of these characters \(\tau\) satisfy (G125.5) in their
unique decomposition \(\tau=\omega_p^r\kappa\).
For every one of them, put \(\psi=\chi_{-4}\tau\).
For both cubic labels and every original finite coefficient
DVR containing these finite character values,
\[
 \operatorname{Sel}_{\rm str}\left(K,
       (E/O)(\chi^{-j}\psi^{-1}c_p^{-e}|_{G_K})\right)=0.
 \tag{G125.14}
\]
This counts characters with conductor \emph{dividing}
\(269^m\), not just characters of primitive conductor
\(269^m\).  It gives no count over other moduli or coefficient primes.
\end{theorem}

\begin{proof}
The Teichmuller subgroup splits the exact reduction map of
\((\mathbb Z/p^m\mathbb Z)^\times\) onto
\(\mathbb F_p^\times\), giving the product of its cyclic
subgroup of order \(268\) and the principal-unit subgroup.
The latter has order \(p^{m-1}\) and is cyclic.  For a
direct check of cyclicity, the binomial expansion gives
\(v_p((1+p)^{p^n}-1)=n+1\) by induction for odd \(p\):
raising a unit \(1+p^a b\), with \(p\nmid b\), to the
\(p\)th power has first term of valuation \(a+1\),
while every later binomial term has larger valuation.
Thus \(1+p\) has order \(p^{m-1}\) modulo \(p^m\).
The finite multiplicative group \(\mathbb F_p^\times\)
is cyclic: its exponent is realized by combining maximal
prime-power orders in the abelian group, and the polynomial
root bound forces that exponent to equal the group order.
These observations also prove the stated cyclic tame factor.

Every finite character therefore has a unique expression
\(\omega_p^r\kappa\), with \(r\) modulo \(268\) and
\(\kappa\) of \(p\)-power order, trivial on the tame factor.
There are \(p^{m-1}\) choices of \(\kappa\), and precisely
two distinct choices
\[
                    r\equiv133-e\quad\text{or}\quad267-e
                                \pmod{268}.
 \tag{G125.15}
\]
This proves (G125.13).  For each character, choose any original
finite coefficient field containing its values.  Reduction of
a root of unity of \(p\)-power order is one, since \(k^\times\)
has order prime to \(p\).  Hence its \(\kappa\) takes values
in \(1+\mathfrak m\).  Cyclotomic extensions of conductor
\(p^m\) are unramified away from \(p\): the roots-of-unity
scheme is finite etale over \(\mathbb Z[1/p]\), since the
derivative of \(X^{p^m}-1\) is invertible at its roots.
Thus its restriction
to \(G_K\) satisfies (G125.2).  The preceding theorem
applies directly and proves (G125.14) over that same \(O\).
A single \(O\) containing all the finitely many character
values may alternatively be fixed before applying the theorem.
The rational characters remain distinct on \(G_K\): the
intersection of \(K\) with \(\mathbb Q(\mu_{p^m})\) is
trivial, since a nontrivial intersection would contain the
quadratic field \(K\), which ramifies at \(3\), whereas the
cyclotomic field does not.  For these finite characters
\(\kappa(-1)=1\) and (G125.5) also gives
\(\psi(-1)=(-1)^{r+1}=(-1)^e\).  This parity observation
is not an assertion of any other original \(\mathrm{GD}\) field.
\end{proof}

\paragraph{Exact consequences and boundary.}
The case \(r=0,\kappa=1\) recovers G114/G115.  The case
\(r=134,\kappa=1\) gives \(\psi=\chi_{-4}\chi_{269}\)
and both weight residues \(133\) and \(267\), using the
quadratic identification already proved in G115.
For every fixed positive weight the finite-character theorem
includes arbitrarily high \(269\)-power conductors and
nonquadratic characters.  Their nonquadratic part may be tame
or wild; the tame part can be quadratic on the original
two weight residues.  These conclusions remove
the corresponding original strict summands.  They do not
supply an automorphic endpoint, model rigidity, a mixed
characteristic-zero Bloch--Kato vanishing, a relaxed local
condition, another coefficient prime, or the complementary
all-family arithmetic budget.  H04.4 and cubic Kummer RH
remain unproved.

\paragraph{Source and geometric dictionary.}
The only non-elementary finite arithmetic inputs are the two
previously proved residual statements: G114 uses the certified
\(D=3\) class bound and its exact nonzero unit image; G115 uses
the certified quotient for \(Y^6-7263\) and its exact unit image.
Their source, certificate and Kummer-descent scopes remain
unchanged.  No fresh certification, complete-unit computation,
Iwasawa control, or automorphic comparison is inferred here.
Hensel lifting, finite-group decomposition and the cohomological
coefficient argument used in this extension are given above.

Geometrically the residual kernel is the same specified
character part of Kummer torsors over the original arithmetic
base used in G114/G115, with trivial class at the coefficient
completion.  Each finite torsion layer of the lifted coefficient
system is a locally constant finite sheaf away from its stated
ramification, by the retained SGA~1, Expose~V, sections~4 and~7
finite-cover dictionary.  The calculation is made over the
original base and passes through the actual uniformizer kernel.
For an infinite-image \(\kappa\), only these finite torsion
layers are asserted to be finite sheaves; their discrete union
is not called a finite cover.  No new motive, polarization or
arithmetic surface is constructed.
